\section*{अध्याय \ref{ch:b3:plant-molecular-physiology} --- आणविक पादप-शरीरक्रिया और प्रतिबल-अनुक्रियाएँ}

\begin{solution}{exo:b3:plant-molecular-physiology:1}
\omterm{def:b3:plant-molecular-physiology:photoreceptors}{फ़ाइटोक्रोम}, लाल \qty{660}{nm} और दूर-लाल \qty{730}{nm}: \omterm{def:b3:plant-molecular-physiology:photoreceptors}{अ-पीतन},
छाया-परिहार, अंकुरण। \omterm{def:b3:plant-molecular-physiology:photoreceptors}{क्रिप्टोक्रोम}, नीला \qty{450}{nm}: \omterm{def:b3:plant-molecular-physiology:photoreceptors}{अ-पीतन}, वह
घड़ी, पुष्पन। \omterm{def:b3:plant-molecular-physiology:photoreceptors}{फ़ोटोट्रोपिन}, नीला: प्रकाश की ओर झुकना, रंध्रों का खुलना,
हरितलवकों का चलना। पर्णहरित प्रकाश को ऊर्जा की तरह नापता है और उस सूचना
से कुछ नहीं करता; प्रकाश-ग्राही रंग, दिशा और अवधि दस लाख गुना कम
फ्लक्सों पर पढ़ते हैं और जीन-अभिव्यक्ति बदल देते हैं --- यानी प्रकाश
खाने और प्रकाश पढ़ने का अंतर।
\end{solution}

\begin{solution}{exo:b3:plant-molecular-physiology:2}
हर \omterm{def:b3:endocrinology:hormone}{हॉर्मोन} किसी दमनकारी के यूबिक्विटिनन तथा विनाश को बढ़ावा देता है:
\omterm{def:b3:plant-molecular-physiology:hormones}{ऑक्सिन} \omterm{def:b3:plant-molecular-physiology:hormones}{TIR1} को Aux/IAA प्रोटीनों से चिपकाता है, GID1 से बँधा \omterm{def:b3:plant-molecular-physiology:hormones}{जिबरेलिन}
\omterm{def:b3:plant-molecular-physiology:hormones}{DELLA} प्रोटीनों को किसी F-बॉक्स लाइगेज़ को सौंप देता है, और जैस्मोनेट
COI1 को JAZ प्रोटीनों से चिपकाता है। वे दमनकारी पहले ही अपने लक्ष्य
जीनों पर बैठे हैं और सक्रियक उनके बगल में तैयार, इसलिए कोई नया प्रोटीन
बनाने की ज़रूरत ही नहीं: उस दमनकारी को नष्ट करना, जो प्रोटिएसोम मिनटों
में कर देता है, उन जीनों को तुरंत चालू कर देता है।
\end{solution}

\begin{solution}{exo:b3:plant-molecular-physiology:3}
ABA PYR/PYL से बँधता है; वह संकुल PP2C फ़ॉस्फ़ेटेज़ का निरोध करता है;
SnRK2 काइनेज़ (OST1), जिसका अब फ़ॉस्फ़ेट नहीं हटता, सक्रिय हो जाती है
और SLAC1 ऋणायन वाहिका का फ़ॉस्फ़ोरिलन करती है; ऋणायन निकलते हैं,
झिल्ली विध्रुवित होती है, बाहरगामी पोटैशियम वाहिकाएँ खुलती हैं,
K$^{+}$ और फिर जल निकलता है, स्फीति गिरती है और वे \omterm{prop:b3:plant-molecular-physiology:stomata}{द्वार कोशिकाएँ} साथ
बैठ जाती हैं। कोई भी कड़ी तोड़ी जा सकती है: कोई ग्राही न हो, कोई ऐसी
फ़ॉस्फ़ेटेज़ हो जिसका निरोध न हो सके (वही चिरप्रतिष्ठित \emph{abi1}
उत्परिवर्ती), OST1 न हो या SLAC1 न हो --- हर एक कोई ऐसा मुरझाता पौधा
देता है जिसके रंध्र सूखे में भी खुले रहते हैं।
\end{solution}

\begin{solution}{exo:b3:plant-molecular-physiology:4}
PTI पूरे सूक्ष्मजीवी वर्गों में साझा अणु (फ्लैजेलिन, काइटिन)
कोशिका-पृष्ठ की ग्राही काइनेज़ों से पहचानती है और किसी मज़बूत भित्ति,
बंद रंध्रों, रक्षा-जीनों की अभिव्यक्ति तथा धीमी वृद्धि में समाप्त होती
है, प्रायः बिना किसी कोशिका-मृत्यु के। ETI किसी विशिष्ट कारक प्रोटीन
को या उसकी क्षति को कोशिका के भीतर बैठे किसी \omterm{def:b3:plant-molecular-physiology:immunity}{NLR ग्राही} से पहचानती है
और संक्रमित कोशिकाओं की अतिसंवेदी मृत्यु तथा किसी तंत्र-व्यापी चेतावनी
में समाप्त होती है।
\end{solution}

\begin{solution}{exo:b3:plant-molecular-physiology:5}
दिन का उजाला $0.87\times 1.15/1.75 = 0.57$; एक पत्ता $0.87\times 0.4/1.0 =
0.35$; सघन वितान $0.87\times 0.1/0.7 = 0.12$; तापदीप्त लैंप
$0.87\times 0.7/1.3 = 0.47$। कोई तंतु-लैंप दूर-लाल से धनी होता है,
इसलिए $\varphi$ दिन के उजाले से नीचे रहता है और वह अंकुर आंशिक छाया
पढ़ता है: वह लंबा हो जाता है। ठंडे सफ़ेद लैंप और बिना दूर-लाल के LED
उलटा करते हैं।
\end{solution}

\begin{solution}{exo:b3:plant-molecular-physiology:6}
भित्ति $5.0$, कोशिकाद्रव्य $7.5$: $(1 + 10^{2.75})/(1 + 10^{0.25}) = 563/2.78
= 203$। भित्ति $4.5$: $563/(1 + 10^{-0.25}) = 563/1.56 = 360$ ---
भित्ति को अम्लीकृत करने से बाहर उदासीन अंश और उसका ग्रहण बढ़ जाता है।
कोशिकाद्रव्य $6.5$, भित्ति $5.5$ के साथ: $(1 + 10^{1.75})/6.62 = 57/6.62 = 8.6$:
फँसाव पाँच गुना गिर जाता है, \omterm{def:b3:plant-molecular-physiology:hormones}{ऑक्सिन} हर फलक से रिसने लगता है, और वे
ध्रुवी प्रवणताएँ चपटी पड़ जाती हैं --- यही एक कारण है कि अवायवीय जड़ें
अपनी दिशा खो देती हैं।
\end{solution}

\begin{solution}{exo:b3:plant-molecular-physiology:7}
\qty{1}{cm/h} $= \qty{10000}{\micro m/h}$: घंटे भर में $200$ कोशिकाएँ,
हर एक में \qty{18}{s}। \qty{50}{\micro m} के आर-पार विसरण: $t =
(5\times 10^{-5})^{2}/(2\times 5\times 10^{-10}) = \qty{2.5}{s}$। उस
कोशिका का भीतरी हिस्सा सेकंडों में पार हो जाता है; उन \qty{18}{s} का
अधिकांश आधारी झिल्ली के PIN से बाहर ढोए जाने की प्रतीक्षा में बीतता है,
और वही दर-सीमक चरण है।
\end{solution}

\begin{solution}{exo:b3:plant-molecular-physiology:8}
$E/A = 1.6\,\Delta w/\Delta c = 1.6\times\num{16000}/150 = 171$ जल-अणु
प्रति CO$_{2}$। C$_{4}$, $\Delta c = 250$: $102$। सूखा दिन,
$\Delta w = 24$: $256$।
\end{solution}

\begin{solution}{exo:b3:plant-molecular-physiology:9}
हिमीकरण कोशिकाओं का जल बाह्यकोशिकीय बर्फ़ में खींचकर उन्हें निर्जलित
कर देता है, इसलिए ठंड और सूखा दोनों निर्जलन-प्रतिबल हैं। वे द्वितीय
संदेशवाहक (ABA, कैल्सियम-स्पंद, क्रियाशील ऑक्सीजन), अनुलेखन कारक
(CBF/DREB कुल दोनों से प्रेरित होता है), और उत्पाद साझा करते हैं:
डिहाइड्रिन तथा LEA प्रोटीन, जो झिल्लियों और प्रोटीनों को जल-हानि से
ढाल देते हैं, संगत विलेय (प्रोलीन, शर्कराएँ), जो जल थामे रखते हैं, और
प्रति-ऑक्सीकारक। जिस पौधे ने उन्हें किसी एक प्रतिबल के लिए बनाया हो
उसके पास वे दूसरे के लिए भी हैं।
\end{solution}

\begin{solution}{exo:b3:plant-molecular-physiology:10}
पूरी धूप के \qty{1}{\%} पर $k_{1} + k_{2} = \qty{0.005}{s^{-1}}$:
साम्य का 95~\% $3/k = \qty{600}{s}$ के बाद, यानी दस मिनट, जबकि दोपहर
में छह सेकंड। आठ घंटे चार अर्ध-आयु हैं: Pfr का $1/16 \approx
\qty{6}{\%}$ बचता है। कोई छोटी कौंध $\varphi$ सेट कर देती है पर फिर वह
Pfr क्षीण हो जाता है, जबकि कोई दिन उसे घंटों ऊँचा रखता है; और जिन
अनुक्रियाओं के लिए Pfr का घंटों काम करना ज़रूरी है (अंकुरण, \omterm{def:b3:plant-molecular-physiology:photoreceptors}{अ-पीतन}) वे
उस संसर्ग का समाकलन करती हैं, इसलिए हल से क्षण भर के लिए ऊपर आया कोई
बीज वैसे अनुक्रिया नहीं देता जैसे सतह पर पड़ा कोई बीज देता है।
\end{solution}

\begin{solution}{exo:b3:plant-molecular-physiology:11}
दिन भर में दस स्थल, हर एक पर $100$ कोशिकाएँ: यानी \num{1000} कोशिकाएँ,
उस पत्ते का $10^{-4}$; और मौसम भर में प्रतिशत का कोई अंश, जबकि कोई एक
संक्रमण बेरोक फैले तो पूरा पत्ता ही जाता। वे मरी कोशिकाएँ लगभग कुछ भी
नहीं लेतीं और उस रोगजनक का भोजन अपने साथ ले जाती हैं --- यानी किसी
जीवपोषी के लिए, जिसे जीवित कोशिकाएँ चाहिए। कोई मृतपोषी मरे ऊतक पर पलता
है: वह \omterm{def:b3:plant-molecular-physiology:immunity}{अतिसंवेदी अनुक्रिया} उसे भोजन और पैर टिकाने की जगह थमा देती है,
और कुछ (\emph{Botrytis}, \emph{Sclerotinia}) जान-बूझकर उसे छेड़ने वाले
विष स्रवित करते हैं; उनके विरुद्ध वह पौधा कोशिका-मृत्यु के बजाय
जैस्मोनेट-मध्यस्थ रक्षाओं पर टिका रहता है।
\end{solution}

\begin{solution}{exo:b3:plant-molecular-physiology:12}
घंटा 12 से 16 तक दर $p$ से उत्पादन। लंबा दिन, भर प्रकाश,
$k = \ln 2/4 = \qty{0.17}{h^{-1}}$: 16 पर स्तर $(p/k)(1 -
\mathrm{e}^{-4k}) = (p/0.17)(0.5) = 2.9p$। छोटा दिन, घंटा 10 से
अँधेरा, $k = \qty{1.39}{h^{-1}}$: $(p/1.39)(1 - \mathrm{e}^{-5.5})
\approx 0.72p$। यानी लंबे दिन में चार गुना अधिक CONSTANS। बीच की कोई
देहली --- मान लीजिए $2p$, जो \emph{FT} सक्रिय करने के लिए चाहिए --- तभी
पार होती है जब उस घड़ी की उत्पादन-खिड़की प्रकाश से संपाती हो: यानी किसी
भीतरी लय का बाहरी दिन से संपात किसी श्रेणीबद्ध दिन-लंबाई को पुष्पन के
किसी सर्व-या-कुछ-नहीं निर्णय में बदल देता है।
\end{solution}

\begin{solution}{pb:b3:plant-molecular-physiology:1}
\textbf{1.} दिन का उजाला $0.87\times 1.15/1.75 = 0.57$; वितान $0.87\times
0.2/0.8 = 0.22$।
\textbf{2.} $r = 2(1 - \varphi)$: दिन के उजाले में \qty{0.86}{mm/h},
छाया में \qty{1.57}{mm/h}; \qty{48}{h} में: $41$ बनाम \qty{75}{mm},
यानी \qty{34}{mm} अधिक।
\textbf{3.} लाल गिरकर दूर-लाल का $0.1\times 1.15 = 0.115$ रह जाता है,
और दूर-लाल गिरकर $0.8$: $\zeta = 0.144$, $\varphi = 0.87\times 0.144/0.744 =
0.17$। हाँ: कोई एक पत्ता $\varphi$ को $0.57$ से $0.17$ पर गिरा देता है,
यानी छाया-परिहार की परास में गहरे।
\textbf{4.} दोनों प्रकाश-परिवर्तन दरें उस फ्लक्स के साथ मापक्रमित होती
हैं, इसलिए उनका अनुपात नहीं। बादल के नीचे वह वर्णक्रम लगभग अपरिवर्तित
रहता है और $\varphi$ $0.57$ के पास ही रहता है: खुले में किसी धुँधले दिन
पर कोई अंकुर लंबा नहीं होता --- वह पड़ोसी पढ़ता है, चमक नहीं।
\textbf{5.} $3/(k_{1} + k_{2})$ के बाद \qty{95}{\%}: पूरी धूप में
\qty{6}{s}, और \qty{1}{\%} पर \qty{600}{s}।
\textbf{6.} \qty{8}{h} $=$ चार अर्ध-आयु, $1/16 = \qty{6.3}{\%}$;
\qty{14}{h} $=$ सात, $1/128 = \qty{0.8}{\%}$। भोर पर बचा Pfr उस रात की
लंबाई बता सकता था, पर वह पलटाव कोई तापीय अभिक्रिया है जो गरम रातों में
तेज़ चलती है, और उसका आरंभिक मान उस दिन के प्रकाश पर निर्भर है: इसलिए
पौधे रात की लंबाई उस घड़ी से नापते हैं।
\textbf{7.} भित्ति: $1/(1 + 10^{0.75}) = 0.15$; कोशिकाद्रव्य: $1/(1 +
10^{2.45}) = 0.0035$।
\textbf{8.} $(1 + 282)/(1 + 5.6) = 283/6.6 = 43$। कोशिकाद्रव्य
$43\times 0.1 = \qty{4.3}{\micro M}$।
\textbf{9.} $\qty{1}{cm/h}/\qty{100}{\micro m} = 100$ कोशिकाएँ प्रति
घंटा, हर एक \qty{36}{s}।
\textbf{10.} \qty{5}{h}। झुकना घंटे भर में शुरू हो जाता है क्योंकि जो
पुनर्वितरण मायने रखता है वह पार्श्व है, यानी शिखर के मिलीमीटर भर के
आर-पार, और वह अनुक्रिया स्थानीय है; वह लंबी दूरी की धारा पूर्ति तय करती
है, समय नहीं।
\textbf{11.} निचला पार्श्व: $60/50 = 1.2$, यानी \omterm{def:b3:plant-molecular-physiology:hormones}{ऑक्सिन} मानक से
\qty{20}{\%} ऊपर, दीर्घन $0.8$; ऊपरी पार्श्व: मानक का $0.8$, दीर्घन
$1.2$। ऊपरी पार्श्व निचले से आगे बढ़ जाता है: वह जड़ नीचे की ओर मुड़ती
है।
\textbf{12.} प्ररोह-कोशिकाएँ अपने ऑक्सिन-इष्टतम से नीचे पड़ी हैं, इसलिए
निचले पार्श्व का अतिरिक्त \omterm{def:b3:plant-molecular-physiology:hormones}{ऑक्सिन} उनकी वृद्धि बढ़ाता है: निचला पार्श्व
ऊपरी से आगे बढ़ जाता है और वह प्ररोह ऊपर की ओर मुड़ता है।
\textbf{13.} $E = 0.3\times 16 = \qty{4.8}{mmol\,m^{-2}\,s^{-1}}$।
\textbf{14.} $A = (0.3/1.6)\times 150 = \qty{28}{\micro mol\,m^{-2}\,s^{-1}}$;
$E/A = 4800/28 = 171$।
\textbf{15.} क्षेत्रफल \qty{0.002}{m^2}, \num{43200}~s: जल $4.8\times
10^{-3}\times 0.002\times\num{43200} = \qty{0.41}{mol} = \qty{7.5}{g}$;
CO$_{2}$ $28\times 10^{-6}\times 0.002\times\num{43200} =
\qty{2.4e-3}{mol}$, यानी $2.4\times 10^{-3}/6\times 180 =
\qty{0.073}{g}$ ग्लूकोस --- यानी ग्राम भर शर्करा के लिए सौ ग्राम जल।
\textbf{16.} $E = 0.05\times 16 = \qty{0.8}{mmol\,m^{-2}\,s^{-1}}$,
$A = \qty{4.7}{\micro mol\,m^{-2}\,s^{-1}}$, अनुपात फिर भी $171$। उस
पौधे ने अपने जल का छह-सातवाँ भाग बचाया और अपने कार्बन का छह-सातवाँ भाग
छोड़ दिया: वाल्व बंद करना दर बदलता है, क़ीमत नहीं।
\textbf{17.} C$_{4}$: $\Delta c = 250$, $E/A = 4800/46.9 = 102$। रात में
CAM: $E = 0.3\times 5 = \qty{1.5}{mmol}$, $A = 28$: $E/A = 53$।
\textbf{18.} $E/A = 1.6\,\Delta w/\Delta c$: वह चालकता कट जाती है
क्योंकि दोनों गैसें वही छिद्र पार करती हैं। वह पौधा उन प्रवणताओं को बदल
सकता है --- भीतर CO$_{2}$ सांद्रित करके (C$_{4}$), हवा नम होने पर
खोलकर (CAM, भोर), पत्ता ठंडा करके, रोओं से सीमा-परत मोटी करके ---
पर उस भौतिकी को नहीं कि दो गैसें एक ही छेद साझा करती हैं।
\textbf{19.} दिन भर में \num{1000} कोशिकाएँ, यानी उस पत्ते का $10^{-4}$।
\textbf{20.} 60 दिनों में \num{60000} कोशिकाएँ, यानी \qty{0.6}{\%}। दिन
में एक बच निकलना, जो \qty{5}{\%} नष्ट करे, उस पत्ते को बीस दिनों में
चुका देता: वह \omterm{def:b3:plant-molecular-physiology:immunity}{अतिसंवेदी अनुक्रिया} उसका सौवाँ भाग लेती है जो कोई एक बचा
हुआ संक्रमण लेता है।
\textbf{21.} प्रकाश-संश्लेषण का $0.2\times 5 = \qty{1}{\%}$। हमेशा चालू
रखने पर वह \qty{5}{\%} लेती और वह पौधा उन पड़ोसियों से पिछड़ जाता जो
उसे पत्तों पर लगाते हैं; प्रेरित रक्षा तभी लाभ देती है जब वह ख़तरा
मौजूद हो।
\textbf{22.} उसे उस NLR के लिए अदृश्यता मिलती है और वह प्रतिरोधी किस्म
को संक्रमित कर लेता है; उसे वह खोना पड़ता है जो वह कारक करता था (PTI
दबाना), इसलिए वह उन पौधों पर दुर्बल है जिनमें वह NLR नहीं। वह
पादप-समष्टि फिर बचे हुए कारकों के विरुद्ध NLR का पक्ष लेने लगती है, और
वह पुराना प्रतिरोध-जीन, जो अब बेकार है, घटता जाता है: यानी दोनों पक्षों
के जीनों का बारंबारता-निर्भर चक्रण।
\textbf{23.} वह नक़ल (कोरोनैटीन) जैस्मोनेट पथ सक्रिय करती है, जो
सैलिसिलेट पथ का विरोध करता है; और सैलिसिलेट ही उन जीवपोषियों के विरुद्ध
रक्षा है जिनमें वह जीवाणु स्वयं है, और वह नक़ल उन रंध्रों को भी फिर खोल
देती है जो उस पौधे ने बंद किए थे। वह रोगजनक उस प्रतिरक्षा-तंत्र की एक
भुजा को दूसरी के विरुद्ध मोड़ देता है।
\textbf{24.} लाखों एक जैसे पौधों पर कोई एक ही प्रतिरोध-जीन कोई एकसमान
वरण है: जो भी उत्परिवर्ती उस कारक को खोए या बदले वह कुछ ही मौसमों में
पूरी फ़सल भर फैल जाता है। प्रजनक कई प्रतिरोध-जीन ढेर करते हैं ताकि कई
कारक एक साथ खोने पड़ें, किस्में मिलाते हैं, वर्षों में जीन बदलते हैं, और
उस व्यापक नमूना-छेड़ी प्रतिरोध का पक्ष लेते हैं जिससे कोई एक उत्परिवर्तन
बच नहीं निकलता।
\textbf{25.} वितान के नीचे $\varphi \approx 0.22$; \omterm{def:b3:plant-molecular-physiology:hormones}{ऑक्सिन} भीतर :
बाहर $\approx 43$; और लगभग $170$ जल-अणु प्रति स्थिर किए गए CO$_{2}$
खोए।
\end{solution}
